\lecture

\begin{guiding}
Everything converges here. The $S_I$ matrix theorem shows that products of
projective spaces are independent in cobordism; the oriented cobordism ring
organizes the question; and Thom's theorem converts cobordism into homotopy
theory. Together they compute $\Omega_*\otimes\Q$.
\end{guiding}

\section{Proof of the matrix theorem}

\begin{proof}[Proof of Theorem~\ref{thm:matrix-nondegenerate}]
\begin{strategy}
Order partitions by their number of parts. The multiplicativity of $S_I$ shows
that $S_I$ can be non-zero on a product indexed by $J$ only when $I$ has at
least as many parts as $J$, and that on the diagonal $I=J$ it is a product of
non-zero numbers. That makes the matrix of the $S_I$ triangular with non-zero
diagonal; a change of basis converts it into the matrix of Chern numbers.
\end{strategy}

\step{1} \emph{Multiplicativity.} Iterating Theorem~\ref{thm:SI-sum} across a
product $X^{j_1}\times\dots\times X^{j_s}$, whose tangent bundle is the sum of
the pull-backs of the $TX^{j_t}$,
\[
S_I\big(X^{j_1}\times\cdots\times X^{j_s}\big)
=\sum_{I_1I_2\cdots I_s=I}
S_{I_1}\big(X^{j_1}\big)\cdots S_{I_s}\big(X^{j_s}\big),
\]
the sum being over all ways of splitting $I$ into $s$ consecutive
sub-partitions.

\step{2} \emph{Which terms survive evaluation.} Evaluating on the fundamental
class, a term vanishes unless each $I_t$ has weight $j_t$, that is, unless
$I_t$ is a partition of $j_t$: otherwise the corresponding factor lies in a
cohomology group above or below the top degree of $X^{j_t}$. In particular each
$I_t$ is non-empty, so
\[
S_I\big[X^{j_1}\times\cdots\times X^{j_s}\big]\neq0
\quad\Longrightarrow\quad
\#\{\text{parts of }I\}\ \ge\ \#\{\text{parts of }J\}.
\]

\step{3} \emph{The diagonal.} If $I=J$ as partitions, the only splittings that
contribute are those with $I_t=(j_t)$ a single part, so
\[
S_J\big[X^{j_1}\times\cdots\times X^{j_s}\big]
=\lambda_J\prod_{t=1}^{s}S_{j_t}\big[X^{j_t}\big],
\]
where $\lambda_J$ is a positive integer counting the number of such splittings.
By hypothesis every factor is non-zero, so the diagonal entry is non-zero.

\step{4} \emph{Triangularity.} Order the $p(n)$ partitions of $n$ so that the
number of parts is non-decreasing. By Step 2 the matrix
$\big(S_I[X^J]\big)_{I,J}$ vanishes below the diagonal blocks and by Step 3 has
non-zero entries on the diagonal, so its determinant is non-zero.

\step{5} \emph{From $S_I$ to Chern numbers.} Each $S_I$ is an integral linear
combination of the monomials $c_{i_1}\cdots c_{i_r}$ of the same weight, and
conversely, because both families are bases of the same lattice, as shown in the
next paragraph. Hence the matrix of Chern numbers is obtained from the matrix of
the $S_I$ by an invertible change of basis and is therefore also
non-degenerate.
\end{proof}

\medskip\noindent\emph{The change of basis.} Inside $\Z[u_1,\dots,u_n]$ the
symmetric polynomials form $\Z[\sigma_1,\dots,\sigma_n]$. Its part of degree $k$
has two natural bases indexed by the partitions $\lambda\vdash k$ into at most
$n$ parts:

\begin{itemize}
\item the monomials in the elementary symmetric functions,
$\sigma_\lambda=\sigma_{\lambda_1}\sigma_{\lambda_2}\cdots$, equivalently
$\sigma_1^{r_1}\cdots\sigma_n^{r_n}$ with $\sum i\,r_i=k$;
\item the symmetrized monomials $S_\lambda$, each distinct monomial of exponent
pattern $\lambda$ appearing once.
\end{itemize}

The first is a basis by the fundamental theorem of symmetric polynomials; the
second because every symmetric polynomial is by definition a sum of
symmetrized monomials. The transition matrix between two bases of the same free
module is invertible over $\Z$, which is what Step 5 used.

\begin{remark}
When applying Theorem~\ref{thm:SI-sum} it matters that each monomial is counted
once. Writing $S_I$ as a sum over all of $S_n$ would count each monomial
$|\mathrm{Iso}|$ times, where $\mathrm{Iso}$ is the isotropy group of the
exponent pattern; the multiplicativity formula holds for the normalized version.
\end{remark}

\medskip\noindent\emph{The Pontryagin version.} For a closed oriented $M^{4m}$
and $I\vdash m$ define $S_I(p(E))$ by the same universal polynomial in the
Pontryagin classes and put $S_I[M]=\langle S_I(p(TM)),[M]\rangle$. Since the
Pontryagin classes are multiplicative only modulo $2$-torsion
(Theorem~\ref{thm:pontryagin-whitney}),
\[
S_I\big(p(E\oplus E')\big)=\sum_{JK=I}S_J\big(p(E)\big)\,S_K\big(p(E')\big)
\qquad\text{modulo $2$-torsion,}
\]
which is enough for evaluation against fundamental classes, torsion classes
pairing to zero with an integral class after multiplication by $2$.

\begin{theorem}\label{thm:pontryagin-matrix}
Let $\{M^{4k}\}_{k\ge1}$ be closed oriented manifolds with $S_k[M^{4k}]\neq0$.
Then for every $n$ the $p(n)\times p(n)$ matrix of Pontryagin numbers
\[
\Big(P_I\big(M^{4j_1}\times\cdots\times M^{4j_s}\big)\Big)_{I,J\vdash n}
\]
is non-degenerate. The family $M^{4k}=\CP^{2k}$ qualifies, since
$p(T\CP^{2k})=(1+u^2)^{2k+1}$ gives $S_k[\CP^{2k}]=2k+1\neq0$.
\end{theorem}

\begin{proof}
The five steps of the proof above apply verbatim, with $c$ replaced by $p$ and
degrees multiplied by two; the only change is that the identities hold modulo
$2$-torsion, which does not affect the evaluations.
\end{proof}

\begin{remark}[quoted]
The classifying space for \emph{oriented} bundles is the oriented Grassmannian
$\tilde G_n$ of oriented $n$-planes, and
\[
H^*\big(\tilde G_{2m+1};\Q\big)=\Q[p_1,\dots,p_m],
\]
\[
H^*\big(\tilde G_{2m};\Q\big)=\Q[p_1,\dots,p_{m-1},e],
\qquad e^2=p_m .
\]
So rationally the Pontryagin classes, together with the Euler class in even
rank, are all the characteristic classes of oriented bundles.
\end{remark}

\section{The oriented cobordism ring}

\begin{definition}
Closed oriented manifolds $M,M'$ of dimension $n$ are \emph{oriented cobordant}
if there is a compact oriented $X$ together with an orientation preserving
diffeomorphism $M\sqcup(-M')\iso\partial X$, the boundary carrying its induced
orientation. Let $\Omega_n$ be the set of equivalence classes.
\end{definition}

\begin{proposition}
$\Omega_n$ is an abelian group under disjoint union, with zero the class of
boundaries and with $-[M]=[-M]$; and
$\Omega_*=\bigoplus_{n\ge0}\Omega_n$ is a graded commutative ring under
Cartesian product.
\end{proposition}

\begin{proof}
Reflexivity uses $M\sqcup(-M)=\partial(M\times[0,1])$; symmetry is reversal of
orientation of the cobordism; transitivity is gluing two cobordisms along a
common boundary component, which is again a smooth manifold after smoothing the
corner. Associativity and commutativity of $\sqcup$ and $\times$ are clear, and
the product of a boundary with a closed manifold is a boundary, so the product
descends to classes.
\end{proof}

\begin{lemma}
Oriented cobordant manifolds have equal Pontryagin numbers.
\end{lemma}

\begin{proof}
If $M\sqcup(-M')$ bounds, Lemma~\ref{lem:pontryagin-boundary} gives
$P_I(M)+P_I(-M')=0$, and $P_I(-M')=-P_I(M')$.
\end{proof}

\begin{corollary}\label{cor:rank-omega}
For every partition $I$ of $k$ the assignment $[M]\mapsto P_I(M)$ is a group
homomorphism $\Omega_{4k}\to\Z$. The products
$\CP^{2i_1}\times\dots\times\CP^{2i_r}$, as $I=(i_1,\dots,i_r)$ runs over the
partitions of $k$, are linearly independent in $\Omega_{4k}$, so
\[
\rk\,\Omega_{4k}\ \ge\ p(k).
\]
\end{corollary}

\begin{proof}
Assemble the homomorphisms into
$\Omega_{4k}\to\Z^{p(k)}$, $[M]\mapsto(P_I(M))_I$. By
Theorem~\ref{thm:pontryagin-matrix} the images of the $p(k)$ products
$\CP^{2i_1}\times\dots\times\CP^{2i_r}$ form a non-degenerate matrix, hence are
linearly independent over $\Q$; so the classes themselves are independent.
\end{proof}

\section{Thom spaces}

\begin{definition}
Let $\pi\colon E\to B$ be a real bundle of rank $k$ with a Euclidean metric, and
$E^{\ge1}=\{e\mid |e|\ge1\}$. The \emph{Thom space} is the pointed space
\[
T(E)=E\big/E^{\ge1},
\qquad
*_E=\big[E^{\ge1}\big],
\]
that is, the unit disc bundle with its boundary sphere bundle collapsed to a
point.
\end{definition}

\begin{center}
\begin{tikzpicture}[scale=0.9]
\draw[thick] (0,0) ellipse (1.1 and 0.34);
\draw[thick] (0,1.5) ellipse (1.1 and 0.34);
\draw[thick] (-1.1,0) -- (-1.1,1.5);
\draw[thick] (1.1,0) -- (1.1,1.5);
\node at (0,-0.85) {\small disc bundle over $B$};
\draw[-{Stealth},thick] (1.9,0.75) -- (3.1,0.75);
\fill (4.6,0.75) circle (1.7pt);
\draw[thick] (4.6,0.75) .. controls (4.05,1.75) and (5.8,1.75) .. (4.6,0.75);
\draw[thick] (4.6,0.75) .. controls (4.05,-0.25) and (5.8,-0.25) .. (4.6,0.75);
\node at (4.7,-0.85) {\small $T(E)$};
\end{tikzpicture}
\end{center}

\begin{proposition}\label{prop:thom-space}
$\tilde H^{i+k}\big(T(E)\big)\iso H^i(B)$ by the Thom isomorphism, for oriented
$E$ over $\Z$. If $B$ is a CW complex then $T(E)$ is a CW complex with one
$(d+k)$-cell for each $d$-cell of $B$, plus the basepoint; in particular
$T(E)$ is $(k-1)$-connected.
\end{proposition}

\begin{proof}
Since $E^{\ge1}$ is a deformation retract of $E_0$, we have
\[
\tilde H^*\big(T(E)\big)\iso H^*\big(E,E^{\ge1}\big)\iso H^*(E,E_0)
\iso H^{*-k}(B)
\]
by Theorem~\ref{thm:thom}. For the cell structure, over a $d$-cell $D^d$ the
bundle is trivial, so the part of the disc bundle above it is
$D^d\times D^k$, contributing one $(d+k)$-cell after the collapse. Since
$T(E)$ then has no cells in dimensions $1,\dots,k-1$ apart from the basepoint, a
cellular approximation argument shows every map $S^m\to T(E)$ with $m<k$ is
null-homotopic.
\end{proof}

\section{Thom's theorem}

\begin{definition}
Let $\mathcal C$ be the class of finite abelian groups. A homomorphism
$h\colon A\to B$ is a \emph{$\mathcal C$-isomorphism} if $\ker h$ and
$\mathrm{coker}\,h$ both lie in $\mathcal C$.
\end{definition}

\begin{theorem}[Hurewicz modulo $\mathcal C$, quoted]
If $X$ is a finite CW complex which is $(k-1)$-connected, then the Hurewicz map
$\pi_r(X)\to H_r(X;\Z)$ is a $\mathcal C$-isomorphism for $r<2k-1$.
\end{theorem}

\begin{corollary}\label{cor:hurewicz-thom}
For an oriented rank $k$ bundle $E$ over a finite CW complex $B$, and $n<k-1$,
the composite
\[
\pi_{n+k}\big(T(E)\big)\longrightarrow H_{n+k}\big(T(E);\Z\big)\iso H_n(B;\Z)
\]
is a $\mathcal C$-isomorphism.
\end{corollary}

\begin{theorem}[Thom]\label{thm:thom-cobordism}
For $k>n+1$ there is a canonical isomorphism
\[
\pi_{n+k}\big(T(\tilde\gamma^k)\big)\ \iso\ \Omega_n,
\]
where $\tilde\gamma^k$ is the universal oriented $k$-plane bundle over
$\tilde G_k$.
\end{theorem}

\begin{proof}[Proof of surjectivity]
\begin{strategy}
Given a closed oriented $M^n$, embed it in a sphere, and collapse everything
outside a tubular neighbourhood; what remains is a map from the sphere to the
Thom space of the normal bundle, which maps to the Thom space of the universal
bundle. The preimage of the zero section is $M$ by construction.
\end{strategy}

\step{1} \emph{Embedding.} By Whitney, embed $M^n\subset\R^{n+k}\subset S^{n+k}$
and let $E$ be its normal bundle, oriented by
$TM\oplus E=T\R^{n+k}|_M$.

\step{2} \emph{The collapse map.} Let $U\iso E$ be a tubular neighbourhood
(Theorem~\ref{thm:tnt}). Collapsing the complement of $U$ gives
\[
S^{n+k}\longrightarrow S^{n+k}\big/\big(S^{n+k}\smallsetminus U\big)\iso T(E),
\]
a continuous map sending everything outside the tube to the basepoint.

\step{3} \emph{To the universal Thom space.} A bundle map $E\to\tilde\gamma^k$,
which exists by the oriented version of Theorem~\ref{thm:universal}, induces a
pointed map $T(E)\to T(\tilde\gamma^k)$. Composing with Step 2 gives
$g\colon S^{n+k}\to T(\tilde\gamma^k)$.

\step{4} \emph{Recovering $M$.} By construction $g$ is smooth and transverse to
the zero section near the preimage of the zero section, and
$g^{-1}(\text{zero section})=M$ with normal bundle $g^*\tilde\gamma^k\iso E$.
Hence the class of $g$ maps to $[M]$.
\end{proof}

\begin{remark}[quoted]
For injectivity, and to define the inverse map, one uses transversality: every
continuous $g\colon S^{n+k}\to T(E)$ is homotopic to one which is smooth and
transverse to the zero section on the preimage of $E^{\le1}$, and then
$g^{-1}(\text{zero section})$ is a closed oriented $n$-manifold whose cobordism
class depends only on the homotopy class of $g$, a homotopy made transverse
providing the cobordism. We do not prove the transversality statements.
\end{remark}

\begin{corollary}
For $k$ large,
\[
\Omega_n\otimes\Q\ \iso\ \pi_{n+k}\big(T(\tilde\gamma^k)\big)\otimes\Q
\ \iso\ H_n\big(\tilde G_k;\Q\big)
=\begin{cases}
\Q^{\,p(m)} & n=4m,\\
0 & \text{otherwise,}
\end{cases}
\]
using Corollary~\ref{cor:hurewicz-thom}, which becomes an isomorphism after
tensoring with $\Q$ since finite groups die, together with the rational
cohomology of $\tilde G_k$ quoted above. Combining with
Corollary~\ref{cor:rank-omega}, where the lower bound $p(m)$ is realized by
products of even complex projective spaces,
\[
\Omega_*\otimes\Q\ =\ \Q\big[\,[\CP^2],\,[\CP^4],\,[\CP^6],\dots\,\big].
\]
\end{corollary}

\begin{remark}[where this leads]
Everything in the course feeds forward. Hirzebruch's signature theorem expresses
the signature of $M^{4k}$ as a universal polynomial in the Pontryagin classes,
precisely because both sides are cobordism invariants and agree on the
generators $\CP^{2i}$. Milnor's exotic spheres are detected by the failure of an
identity between Pontryagin numbers. And index theory reinterprets all of these
numbers analytically.
\end{remark}
